%% -*- coding: utf-8 -*-
\documentclass[12pt,pagesize,paper=192mm:108mm,landscape]{scrbook} 
%1920x1080 1280x720
\areaset[current]{192mm}{108mm}
\usepackage{calc}
\usepackage[T2A]{fontenc}
\usepackage[utf8]{inputenc}
\usepackage[english,russian]{babel}
\usepackage{microtype}
\usepackage{misccorr}
\usepackage{cmap}
%\usepackage[unicode=true]{hyperref}
\usepackage{graphicx}
\usepackage{amssymb}
\usepackage{amsmath}
%\usepackage{srcltx}
\usepackage{textcomp}
\usepackage{xspace}
%научные символы и смайлики \smiley \frownie
\usepackage{wasysym}
\usepackage{ccicons}
%перенос формул в тексте
\newcommand*{\hm}[1]{#1\nobreak\discretionary{}%
  {\hbox{$\mathsurround=0pt #1$}}{}}

\begin{document}
\begin{titlepage}
  \vspace*{-0.5em}
  \begin{center}    
    % \hspace*{3em}
    % \begin{minipage}[t]{3em}
    %   \includegraphics[width=\textwidth]{../BINP-logo}
    % \end{minipage}\hfill
    % \begin{minipage}{0.23\linewidth}
    % \includegraphics[width=\textwidth]{../NSU-logo}
    % \end{minipage}
    % \hfill
    % \hspace*{6em}

    % Кафедра теоретической физики физического факультета НГУ
    % \medskip

    % \Large
    % Профессор Грабовский А.\,В.
    % \smallskip

    % \huge
    % \textbf{Общая теория относительности}
    % \smallskip

    \vfill
    \Large
    Задание № 9
 \smallskip

    \normalsize
    \begin{minipage}{0.9\linewidth}
      \begin{enumerate}
      \item Показать, что плотность энергии гравитационной волны
        выражается через метрику волн с~``+'' и ``х'' поляризациями так
        \[t_{00} = \frac{c^4}{16\pi G}\cdot \left<\left(h_{+,t}\right)^2+\left(h_{x,t}\right)^2\right>.\]
      \item Показать, что скорость изменения периода обращения двойной
        системы с разными массами получается из формулы для системы с
        одинаковыми массами заменой 
        \[(2M)^{\frac{5}{3}} \to 4M_1M_2/\left(M_1+M_2\right)^{\frac{1}{3}}.\]
      \item Показать, что скорость изменения периода обращения двойной
        системы по орбите с~ненулевым эксцентриситетом $e$ отличается от
        формулы для нулевого эксцентриситета множителем
        \[\left(1 + \frac{73}{24}\cdot e^2 + \frac{37}{96}\cdot e^4\right)\big/\left(1-e^2\right)^\frac{7}{2}.\]
      \end{enumerate}
    \end{minipage}
    \vfill

    % \normalsize \ccbysa\hspace{0.5em}  Новосибирск 2022
  \end{center}
\end{titlepage}
\end{document}
